\documentclass[11pt]{article}
\usepackage{preprint}

\title{Data Collection and Accountability}
\author{Jiadi He\email{jh2288@georgetown.edu}}
\shorttitle{Data Collection and Accountability}
\keywords{data; measurement; institutional accountability}

\begin{document}
\maketitle
\keywordsline

\section{Introduction}

When there is a need to improve systems or programs, the first stage is to collect data to help identify the areas that need improvement. Data informs the developers of policies and those tasked with designing new programs in the areas where the old systems fell short and what needs the new programs must address. The collection of data is thus focused on the identified needs. The researchers believe that data will provide an excellent basis for finding the best solution. However, that might not be the case. Data may provide a basis for manipulating the present systems to attain desired results without necessarily going through a sustainable process. For instance, when the quality education test passes, teachers focus on teaching the exams while ignoring other important areas of learning that are not likely to get tested. There is an argument which advance that the contemporary data collection methods are insufficient and inadequate in painting the real picture of the level and quality of education in different schools

\section{Somerville School Framework}

Jack Schneider and his team were tasked with improving the school quality in Somerville. They decided to collect new data instead of using the available data because, in their previous experience, the use of available data had led to significant dissatisfaction. Data collection would engage stakeholders, teachers, parents, the public and policymakers (Schneider, 2017). The polls sought to identify five major things that could be adopted to make regional schooling perform well. The investigation yielded varied answers. Surprisingly, there were similarities among some groups on what was necessary for a good school. Eventually, the results compiled and sent to the district for approval did not vary greatly as the investigators condensed the answers to what they thought would provide a simplified observation of needs from the responses. The results were categorized into essential inputs. The first input entailed what was referred to as teachers and the teaching environment. The second category was the school culture, and the third metric was resources. These were further broken down into subcategories to give specific inputs the schools should adopt (Schneider, 2017). The findings made it possible for schools to put specific measures as basic minimum requirements for performance.

Schneider (2017) argues that standardized tests do not determine the quality of learning as they are insufficient in capturing important elements that do define education quality. He argues that the data provided by tests on the quality of education in institutions do not provide the true picture of quality in different institutions. He argues that the demographics of schools are more telling of the quality of learning than standardized tests. His investigation is led by the principle that a quality learning environment will contribute to learning. Passing the standardized test will be one of the indications that learning has occurred. However, his position is not supported by the arguments advanced by Campbell's Law which advances that when institutions seek to optimize metrics without considering other equally important factors the end up with the corruption of those metrics. Schneider goes against this law by creating a list of important metrics that organization must observe to improve the quality of education (Asch \& Musgrove, 2017). The data collection and accountability advanced by Schneider (2017) seek to identify certain metrics that must be present in learning. He argues for a replacement of standardized tests with several other metrics designed to improve the quality of the environment in which learning occurs. The quantitative social indicators he suggests should be used for decision making are subject to corruption and distortion.

Schneider (2017) and Koretz (2008) appear to agree on how quality education should be fundamentally designed. The two scholars agree that the use of test data generated in school is not sufficient to determine the quality of learning. Koretz argues that tests can only be a subset of the goals of education and that they generally provide for a small sample of behavior used to estimate students' mastery over extensive domains of knowledge and skills. Tests can also not provide the best data source since they provide inconsistent results. Koretz (2008) also mentions that tests do not provide a consistent manner in which the results can be reported. The arguments are consistent with those raised by Schneider. However, Koretz appears to accept that testing might be necessary for gauging the competence (while it is not 100\% foolproof) of some of the important areas of learning in key sectors such as medicine, psychology and learning. He thus does not call tests unnecessary in the same way Schneider does. He believes they are necessary but not efficient. What Koretz and Schneidder fail to include in their positions is the issue of accountability. While they argue for the adoption of different metrics, they do not provide a framework that will ensure that the new metrics will be accountable and not subject to corruption as is expected under Campbell’s law.

\section{Conclusion}

There is a significant need to ensure that the process adopted does not corrupt the system it is meant to improve by subverting the process to attain the expected outcomes without emphasizing the importance of a well-thought process. His argument drives the position taken by Schneider that the quality of learning should not be determined by standardized tests. He argues that the data provided by tests on the quality of education in institutions do not provide the true picture of quality in different institutions. He provides alternative metrics and thus goes against the position advanced in Campbell's law. However, Schneider and Koretz appear to agree on how quality education should be fundamentally designed. The two scholars agree that the use of test data generated in school should not be sufficient to determine the quality of learning.

\begin{references}
\item Asch, C. M., \& Musgrove, G. D. (2017). \textit{Chocolate City: A History of Race and Democracy in the Nation's Capital}. UNC Press Books.
\item Koretz, D. (2008). \textit{Measuring up: What educational testing really tells us}. Harvard University Press.
\item Schneider, J. (2017). \textit{Beyond test scores: A better way to measure school quality}. Harvard University Press.
\end{references}

\end{document}
